% GENERATED FILE. Edit canonical lessons, then run npm run generate:books.
% canonical-source-hash: fnv1a64:ca871c010650a92d
% canonical-lessons: FA-W288-ba-saat

\chapter{Writing --- With A Clock}
\label{ch:fa-with-a-clock}
\hlchaptermodality{288}

\begin{chapteropening}
\textbf{By the end of this chapter:} \emph{I can fill in a short practical form in Persian script against a three-minute clock, and I know what this track's A1 writing paper asks for.}

Complete FA-W288-ba-saat: all six form fields in Persian script before the three-minute timer rings, then stop.
\end{chapteropening}

\section[(short timed A1 writing checkpoint)]{The same writing, with a clock}

\label{lesson:FA-W288-ba-saat}



\begin{quote}
You wrote your own line with as long as you liked. A task you can only do untimed is not yet a task you can do.

This page will not show you any copyable Persian answer. Put away the earlier lessons, romanization, dictionaries, translators, and spell-checkers.
\end{quote}

\subsection*{You'll want to know — what the A1 paper is}
This book found no widely used international A1 writing paper for Persian. \textbf{SAMFA}, the Iranian exam this book names as its far destination, is one academic test: its writing section is two essays in one hour, and the provider notice this book checked maps it to no CEFR level at all. So there is no SAMFA A1 paper to rehearse.

This book's A1 writing paper is therefore its own, modelled on what the Council of Europe's framework expects of A1 writing: a form with personal details, and a short, simple message. It runs \textbf{20 minutes}: a practical form, then a \textbf{30 to 40 word} Persian message to a named reader. Today you rehearse the form, for three minutes.

\subsection*{Writing — timed assessment production}
\textbf{One practical form — three minutes}

Set a \textbf{three-minute timer}. Complete \textbf{all six fields in Persian script} for a new neighbourhood reading group. Write one short answer in every field, right to left. Start the timer when you begin the first field; leave any unfinished field as evidence when it rings.

\begin{enumerate}
\raggedright
  \item Your name
  \item How you are today
  \item One person in your family
  \item One drink you would like at the meeting
  \item One thing you will bring
  \item A goodbye
\end{enumerate}

When the three-minute timer rings, \textbf{stop}. Record the number of completed fields. An honest short attempt is more useful than a perfect untimed one.

\begin{tcolorbox}[breakable,colback=teal!4,colframe=teal!35!black,title={Before you move on}]
Spend at most one minute looking at this first attempt. Check the dots and the joins in one field, note one field to practise next, and do not rewrite the timed answer. Keep the page as evidence.

One fact about the full assessment changes how you study. Its four papers are each out of 100, and you need \textbf{60 on every one} in the same sitting. Nothing averages: strong reading cannot carry weak writing. And one of writing's four scored dimensions is Persian orthographic control, so a missing dot costs marks.

This checkpoint is \textbf{not the complete A1 writing paper}, which still needs its mocks, rubric, calibration, and human review. It does not establish A1 attainment, readiness, completeness, or human validation.
\end{tcolorbox}
